\chapter{Conception}

\section{Introduction}

Dans le contexte de ce projet, nous accordons une importance
particulière à la conception en raison de la nature hybride du système à produire~:
une application web Django intégrant un service de prédiction externe
FastAPI, ainsi qu'un pipeline de \textit{machine learning} reposant sur le
modèle Random Forest que nous avons sélectionné au chapitre précédent. Nous avons veillé à ce que ces trois composants interagissent de manière fluide
tout en restant suffisamment découplés pour évoluer indépendamment.

Dans ce chapitre, nous présentons successivement l'architecture globale du système
et les motivations de nos choix, la conception détaillée de
l'API d'inférence FastAPI, la conception de la base de données, la
modélisation des classes métier, les diagrammes de séquence des flux
principaux, la conception de l'interface utilisateur, et enfin nous
justifions les choix technologiques que nous avons retenus.
\section{Planification des sprints applicatifs}

Conformément à la méthodologie Scrum que nous avons adoptée, nous nous concentrons dans cette phase sur les sprints dédiés à la réalisation de la plateforme web et de l'API d'inférence. Le tableau suivant synthétise les objectifs opérationnels et les flux techniques que nous couvrons dans ce cycle.

\begin{table}[htbp]
\centering
\caption{Synthèse des sprints du volet applicatif}
\label{tab:recap_sprints_applicatifs}
\begin{tabular}{|c|p{8cm}|c|}
\hline
Sprint & Objectif principal et fonctionnalités clés & Total SP \\
\hline
S4 & \textbf{API et Sécurité :} Conception de l'API FastAPI (Lifespan, Pydantic), flux d'authentification 2FA avec OTP, et système de contrôle d'accès RBAC. & 51 \\
\hline
S5 & \textbf{Intégration et Données :} Flux d'import multi-formats, moteur d'isolation par agence (middleware), et flux de prédiction batch avec mécanisme de failover local. & 34 \\
\hline
S6 & \textbf{Métier et Interface :} Tableau de bord dynamique par rôle (Chef vs Agent), flux de validation des recommandations, notifications in-app et export PDF. & 40 \\
\hline
Total & & 125 \\
\hline
\end{tabular}
\end{table}

Cette planification nous permet de structurer les flux d'interaction entre les composants, en garantissant que chaque fonctionnalité critique (comme l'isolation des agences ou la sécurité 2FA) est intégrée dès la phase de conception technique.

\section{Architecture générale du système}

\subsection{Vue d'ensemble}

Nous avons opté pour une architecture hybride à trois niveaux enrichie
d'un micro-service d'inférence ML externe. Cette organisation nous permet de découpler
explicitement la logique métier, la persistance des données et le
traitement ML en trois couches aux responsabilités bien délimitées.

\begin{figure}[htbp]
    \centering
    \includegraphics[width=0.7\linewidth,height=7cm]{images/architecture globale.png}
    \caption{Architecture globale du système}
    \label{fig:architecture_globale}
\end{figure}

Le flux de traitement que nous avons mis en place suit le cheminement suivant~: le
navigateur de l'utilisateur émet des requêtes HTTP vers l'application
Django, qui s'exécute sur le port~8000. Django orchestre l'ensemble
de la logique métier, persiste les données dans la base PostgreSQL
via son ORM, et délègue les opérations d'inférence ML au service
FastAPI qui tourne sur le port~8001 (interne 8000). Nous avons intégré un pipeline ML local qui prend automatiquement le relais en cas d'indisponibilité de ce
dernier, contribuant ainsi à la continuité de service sans intervention de l'utilisateur.

\subsection{Justification de nos choix architecturaux}
\begin{enumerate}
\item Séparation Django et FastAPI

Nous avons fait le choix d'isoler le moteur de prédiction ML dans un micro-service FastAPI
indépendant plutôt que de l'intégrer directement dans Django. Ce choix
repose sur trois motivations complémentaires que nous détaillons ici.

D'un point de vue technique, nous avons privilégié FastAPI car il est construit sur ASGI
(\textit{Asynchronous Server Gateway Interface}) et offre des
performances nativement asynchrones particulièrement adaptées aux
opérations de prédiction batch. Django, dont l'architecture est
historiquement synchrone, aurait bloqué le processus serveur pendant
toute la durée du calcul.

D'un point de vue de maintenabilité, cette isolation nous permet
de faire évoluer, remplacer ou réentraîner le modèle de prédiction
sans impacter l'application Django. Nous avons formalisé le contrat d'interface entre les
deux services par des schémas Pydantic définis dans le
service FastAPI (\texttt{churn\_api}), constituant ainsi une frontière
explicite et vérifiable.

Enfin, cette architecture nous ouvre la voie à un déploiement
différencié~: nous pouvons héberger le service FastAPI sur un serveur
disposant de ressources matérielles adaptées aux calculs ML, tandis
que Django reste sur une infrastructure web standard.

\item Organisation modulaire de l'application Django

Nous avons décomposé l'application Django en quatre modules autonomes,
chacun portant un périmètre fonctionnel clairement défini.

\begin{table}[htbp]
    \centering
    \caption{Responsabilités des modules Django}
    \label{tab:modules_django}
    \begin{tabular}{|p{2.8cm}|p{4cm}|p{6.5cm}|}
        \hline
        \textbf{Module} & \textbf{Couche applicative}
        & \textbf{Responsabilité principale} \\
        \hline
        \texttt{accounts}
        & Authentification
        & Cycle de vie des comptes utilisateurs, authentification à
          double facteur (OTP), workflow d'inscription et de
          validation hiérarchique, journalisation des connexions. \\
        \hline
        \texttt{core}
        & Services partagés
        & Service OTP centralisé, moteur de notifications, client HTTP vers FastAPI,
          modèles géographiques \texttt{Ville}/\texttt{Agence},
          \textit{middleware} et \textit{context processors}. \\
        \hline
        \texttt{dashboard}
        & Interface et métier
        & Tableau de bord, gestion des recommandations et de leur
          workflow de validation, notifications in-app, historique
          des analyses, export PDF, réinitialisation du dashboard. \\
        \hline
    \end{tabular}
\end{table}
\begin{table}[htbp]
    \centering
    \begin{tabular}{|p{2.8cm}|p{4cm}|p{6.5cm}|}
        \hline
        \texttt{learning}
        & Données clients
        & Modèles \texttt{Dataset} et \texttt{ClientChurn}, import
          CSV, générateur de données mock, commande de gestion
          \texttt{import\_dataset}. \\
        \hline
    \end{tabular}
\end{table}
\end{enumerate}

\subsection{Conception du système de recommandation stratégique}

Nous avons conçu le moteur de recommandation non comme un simple outil technique, mais comme le bras armé de la stratégie marketing de l'entreprise. Sa conception repose sur quatre principes fondamentaux du marketing relationnel moderne :

\begin{enumerate}
    \item \textbf{La Rétention Proactive :} Plutôt que de réagir au départ du client, le système anticipe le churn. Nous avons modélisé un workflow où chaque client à risque est automatiquement transformé en une mission opérationnelle pour un agent, garantissant qu'aucun client fragile ne soit ignoré.
    
    \item \textbf{L'Hyper-Personnalisation (Explainable AI) :} Grâce à l'intégration des valeurs SHAP, chaque recommandation est justifiée par les comportements réels du client (ex: baisse de data, réclamations). Cette conception permet à l'agent d'avoir une approche commerciale "chirurgicale" et non générique.
    
    \item \textbf{La Priorisation par la Valeur (CLV) :} Dans un contexte de ressources limitées, nous avons intégré le calcul de la \textit{Customer Lifetime Value}. Le système oriente les efforts de l'agence vers les clients représentant le plus fort enjeu financier, optimisant ainsi le ROI des campagnes de rétention. La formule simplifiée de la CLV utilisée dans notre système est définie par :
    
    \begin{equation}
        CLV = ARPU \times \left( \sum_{t=1}^{H} (1 - P_{churn}) \right) \approx ARPU \times \lfloor (1 - P_{churn}) \times H \rceil
    \end{equation}
    
    Où :
    \begin{itemize}
        \item $ARPU$ (\textit{Average Revenue Per User}) est représenté par la facture moyenne mensuelle du client.
        \item $P_{churn}$ est la probabilité de churn prédite par le modèle ML.
        \item $H$ est l'horizon temporel de rétention (fixé à 12 mois dans notre application).
    \end{itemize}
    
    Les recommandations sont ensuite triées selon un score de priorité multicritère $S_{p}$ défini par :
    
    \begin{equation}
        S_{p} = \min(P_{metier}) + \frac{1}{CLV}
    \end{equation}
    
    Ce qui permet de traiter en priorité les règles ayant la plus haute importance métier ($P_{metier}$ le plus bas), tout en utilisant la $CLV$ comme facteur de départage en cas d'égalité.
    
    \item \textbf{La Responsabilité Opérationnelle (Closed-Loop) :} Nous avons conçu un circuit de validation hiérarchique. Une recommandation suit un cycle de vie strict : Proposition $\rightarrow$ Validation Chef $\rightarrow$ Exécution Agent $\rightarrow$ Clôture. Ce principe assure une traçabilité complète de l'effort commercial.
\end{enumerate}

\subsection{Conception détaillée de l'API d'inférence FastAPI}

\begin{enumerate}
    \item Choix du framework

Nous avons choisi FastAPI pour trois raisons déterminantes.
Premièrement, pour ses performances~: il est l'un des frameworks Python les plus rapides. Deuxièmement, pour sa
validation automatique via Pydantic qui nous garantit l'intégrité des données sans code supplémentaire. Troisièmement, nous avons apprécié sa
documentation automatique via Swagger UI, qui a facilité nos phases de tests et
d'intégration.

\item Gestion du cycle de vie~: pattern Lifespan

Nous initions le service FastAPI dans \texttt{main.py} via le
pattern Lifespan. Nous avons configuré le système pour que les artefacts lourds soient chargés \textit{une
seule fois} au lancement du serveur et stockés dans un dictionnaire
global. Ce
mécanisme nous garantit l'accessibilité des données par tous les endpoints et l'intégrité du service au démarrage.

\item Validation des données~: Pydantic

Dans le module \texttt{schemas/}, nous avons défini le schéma \texttt{ClientInput} pour
garantir l'intégrité de chaque requête entrante. Pydantic vérifie
automatiquement pour nous les types et les contraintes logiques~:

\begin{itemize}
    \item \texttt{Field(..., ge=0)}~: nous imposons que la valeur soit positive ou nulle, empêchant
    des données absurdes d'atteindre le modèle.

    \item \texttt{Field(..., ge=1, le=5)}~: pour
    \texttt{satisfaction\_client}, nous avons contraint la plage à
    l'échelle de l'enquête originale (1 à~5).
\end{itemize}

Nous avons conçu le schéma de sortie \texttt{PredictionOutput} pour retourner un objet
structuré contenant la probabilité de churn, la décision binaire
basée sur notre seuil métier de 0,32, et les explications SHAP.

\item Architecture du projet et artefacts ML

Nous avons organisé l'architecture de l'API selon une séparation stricte
entre les artefacts de modélisation et le code de production~:

\begin{figure}[htbp]
\centering
\begin{lstlisting}[basicstyle=\ttfamily\small]
churn_api/
|
|-- Dockerfile
|-- requirements.txt
|
|-- fastapi_artifacts/
|   |-- churn_model_v1.pkl       <- Modele serialise (Random Forest)
|   |-- churn_threshold_v1.pkl   <- Seuil optimal (0,32)
|   |-- feature_names_v1.json    <- Ordre des features (critique)
|   |-- churn_metadata_v1.json   <- Metriques et metadonnees
|   |-- shap_explainer_v1.pkl    <- Explainer TreeSHAP
|   `-- shap_meta_v1.json        <- Metadonnees SHAP
|
`-- app/
    |-- main.py                  <- Point d'entree de l'application
    |-- routers/                 <- Definition des endpoints
    |-- schemas/                 <- Schemas Pydantic (I/O)
    `-- ml/                      <- Coeur du moteur d'inference
\end{lstlisting}
\caption{Structure du projet backend (churn\_api)}
\label{fig:archi_projet}
\end{figure}

Nous soulignons trois décisions d'architecture majeures. L'isolement des artefacts dans \texttt{fastapi\_artifacts/} constitue le pont entre la science des
données et l'ingénierie logicielle. Le versionnage des artefacts
(\texttt{\_v1}) nous permet de déployer une version améliorée du modèle
sans interrompre le service existant. Enfin, nous considérons le fichier
\texttt{feature\_names\_v1.json} comme le point le plus critique~: sans l'ordre exact des colonnes, le modèle
produirait des prédictions erronées en production.

\subsubsection{Description détaillée des endpoints}

Nous avons exposé plusieurs endpoints regroupés en trois familles
fonctionnelles.

\paragraph{\texttt{GET /health}~ monitoring du système.}
Nous utilisons cet endpoint pour vérifier que l'API est active et que
les artefacts sont chargés. Django le sollicite avant chaque analyse.

\paragraph{\texttt{GET /api/model/info}~ gouvernance ML.}
Nous y exposons les métadonnées du modèle (version, Recall, seuil optimal) pour permettre un audit des performances en production.

\paragraph{\texttt{POST /api/predict}~ prédiction individuelle.}
C'est notre endpoint central. Nous y avons implémenté un flux en trois étapes~: l'alignement des features, la décision basée sur le seuil de 0,32, et le calcul des valeurs SHAP pour l'interprétabilité locale.

\paragraph{\texttt{POST /api/predict/batch}~ scoring par lot.}
Nous avons conçu cet endpoint pour traiter une liste de clients en une seule
requête, produisant ainsi une liste de priorités d'intervention pour les équipes.
\end{enumerate}
\section{Conception de la base de données}

\subsection{Système de gestion de base de données retenu}

Nous avons retenu PostgreSQL comme système de gestion de base de données. Ce choix nous assure une
homogénéité entre le développement et la production. Nous avons veillé à externaliser les paramètres de connexion dans le fichier
\texttt{.env} pour des raisons de sécurité.

\subsection{Dictionnaire de données}

Afin de refléter la complexité réelle du système et d'assurer une traçabilité complète des informations, nous avons élaboré un dictionnaire de données exhaustif. Il regroupe l'intégralité des tables organisées par modules fonctionnels.

\begin{footnotesize}
\begin{longtable}{|p{2.5cm}|p{4.5cm}|p{6cm}|}
    \caption{Dictionnaire exhaustif des données}
    \label{tab:dict_donnees} \\
    \hline
    \textbf{Entité} & \textbf{Description} & \textbf{Attributs principaux (Type)} \\
    \hline
    \endfirsthead
    \hline
    \textbf{Entité} & \textbf{Description} & \textbf{Attributs principaux (Type)} \\
    \hline
    \endhead
    \hline
    \endfoot
    
    \multicolumn{3}{|l|}{\textit{Module de gestion des accès (Accounts \& Core)}} \\
    \hline
    \texttt{Ville} & Référentiel des villes . & \begin{itemize}[leftmargin=*, nosep] \item \texttt{nom} (String) \item \texttt{region} (String) \item \texttt{prioritaire} (Bool) \end{itemize} \\
    \hline
    \texttt{Agence} & Points de vente rattachés à une ville. & \begin{itemize}[leftmargin=*, nosep] \item \texttt{nom} (String) \item \texttt{ville} (FK) \item \texttt{code} (String) \end{itemize} \\
    \hline
    \texttt{User} & Utilisateurs avec rôles et habilitations. & \begin{itemize}[leftmargin=*, nosep] \item \texttt{username} (String) \item \texttt{role} (Enum) \item \texttt{agence} (FK) \item \texttt{statut} (Enum) \end{itemize} \\
    \hline
    \texttt{OTPCode} & Codes de sécurité à usage unique. & \begin{itemize}[leftmargin=*, nosep] \item \texttt{user} (FK) \item \texttt{code} (String) \item \texttt{expire\_at} (DateTime) \end{itemize} \\
    \hline
    \texttt{AdminVille} & Supervision régionale des administrateurs. & \begin{itemize}[leftmargin=*, nosep] \item \texttt{user} (FK) \item \texttt{ville} (FK) \item \texttt{actif} (Bool) \end{itemize} \\
    \hline
    \texttt{LoginActivity} & Journal d'audit des connexions. & \begin{itemize}[leftmargin=*, nosep] \item \texttt{user} (FK) \item \texttt{timestamp} (DateTime) \end{itemize} \\
    \hline

    \multicolumn{3}{|l|}{\textit{Module de données et Intelligence Artificielle (Learning)}} \\
    \hline
    \texttt{Dataset} & Source et métadonnées d'un import. & \begin{itemize}[leftmargin=*, nosep] \item \texttt{nom} (String) \item \texttt{methode} (Enum) \item \texttt{nb\_clients} (Int) \end{itemize} \\
    \hline
    \texttt{ClientChurn} & Profil complet et scores prédictifs. & \begin{itemize}[leftmargin=*, nosep] \item \texttt{client\_id} (String) \item \texttt{score\_churn} (Float) \item \texttt{churn\_predit} (Bool) \end{itemize} \\
    \hline
    \texttt{EvenementCDR} & Historique des appels et consommations. & \begin{itemize}[leftmargin=*, nosep] \item \texttt{client} (FK) \item \texttt{type\_evenement} (Enum) \item \texttt{duree\_appel} (Int) \end{itemize} \\
    \hline
    \texttt{ShapValeur} & Décomposition SHAP de la prédiction. & \begin{itemize}[leftmargin=*, nosep] \item \texttt{client} (FK) \item \texttt{feature\_name} (String) \item \texttt{value} (Float) \end{itemize} \\
    \hline

    \multicolumn{3}{|l|}{\textit{Module opérationnel et pilotage (Dashboard)}} \\
    \hline
    \texttt{AnalyseSession} & Résumé et métriques d'une analyse. & \begin{itemize}[leftmargin=*, nosep] \item \texttt{taux\_churn} (Float) \item \texttt{recall} (Float) \item \texttt{date} (DateTime) \end{itemize} \\
    \hline
    \texttt{Recommandation} & Actions de rétention assignées. & \begin{itemize}[leftmargin=*, nosep] \item \texttt{statut} (Enum) \item \texttt{type} (Enum) \item \texttt{assignee\_a} (FK) \end{itemize} \\
    \hline
    \texttt{RejetRec} & Workflow de rejet des actions. & \begin{itemize}[leftmargin=*, nosep] \item \texttt{recommandation} (FK) \item \texttt{motif} (Text) \item \texttt{statut} (Enum) \end{itemize} \\
    \hline
    \texttt{Notification} & Alertes in-app pour les agents. & \begin{itemize}[leftmargin=*, nosep] \item \texttt{destinataire} (FK) \item \texttt{titre} (String) \item \texttt{lu} (Bool) \end{itemize} \\
    \hline
    \texttt{Message} & Communication interne sur fiches clients. & \begin{itemize}[leftmargin=*, nosep] \item \texttt{destinataire} (FK) \item \texttt{sujet} (String) \item \texttt{lu} (Bool) \end{itemize} \\
    \hline
\end{longtable}
\end{footnotesize}
  
\section{Diagrammes de classes}

Dans cette section, nous décrivons la structure statique de notre système. Nos classes métier correspondent
directement aux modèles Django que nous avons définis.

\subsection{Module \texttt{accounts}}

Nous y regroupons la gestion des utilisateurs, les codes OTP et le journal d'audit des connexions.

\begin{figure}[htbp]
    \centering
    \includegraphics[width=0.6\linewidth,height=9cm]{images/classacc.png}
    \caption{Diagramme de classes du module \texttt{accounts}}
    \label{fig:dc_accounts}
\end{figure}

\subsection{Module \texttt{learning}}

Ce module concentre les entités liées aux données clients et aux résultats ML que nous avons conçues.

\begin{figure}[htbp]
    \centering
    \includegraphics[width=0.7\linewidth, height=9cm]{images/classlearn.png}
    \caption{Diagramme de classes du module \texttt{learning}}
    \label{fig:dc_learning}
\end{figure}

\subsection{Module \texttt{dashboard}}

Nous y regroupons la logique liée aux recommandations, aux notifications et au suivi des sessions d'analyse.

\begin{figure}[htbp]
    \centering
    \includegraphics[width=0.7\linewidth, height=7.5cm]{images/classdash.png}
    \caption{Diagramme de classes du module \texttt{dashboard}}
    \label{fig:dc_dashboard}
\end{figure}

\section{Diagrammes de séquence}

Nous utilisons les diagrammes de séquence pour modéliser les interactions dynamiques entre les composants de notre architecture. Ces schémas illustrent comment les différents services (Serveur d'Application, Moteur d'Inférence ML, PostgreSQL, SendGrid) collaborent pour réaliser les processus métier critiques, en incluant la gestion des scénarios alternatifs.

\subsection{Flux 1~: Authentification à double facteur (2FA)}

Pour sécuriser l'accès à la plateforme, nous avons implémenté un flux d'authentification robuste. Ce processus commence lorsque l'utilisateur soumet ses identifiants au Serveur d'Application. Nous utilisons un fragment combiné \textit{alt} pour gérer la première vérification : si les identifiants sont invalides, l'accès est immédiatement refusé ; s'ils sont corrects, nous déclenchons la génération d'un code OTP.

Le Serveur d'Application interagit ensuite avec l'API SendGrid pour acheminer ce code. La session reste dans un état intermédiaire jusqu'à la saisie de l'utilisateur. Un second fragment \textit{alt} traite la validation de l'OTP : en cas d'expiration ou de code erroné, la session est invalidée ; en cas de succès, nous ouvrons la session finale et enregistrons l'activité de connexion.

\begin{figure}[htbp]
    \centering
    \includegraphics[width=0.9\linewidth]{images/auth.png}
    \caption{Diagramme de séquence de l'authentification à double
    facteur avec fragments alternatifs}
    \label{fig:ds_authentification_2fa}
\end{figure}

\subsection{Flux 2~: Analyse de portefeuille avec FastAPI}

Ce flux constitue le coeur technologique de notre solution. Lorsqu'un Chef d'Agence lance une analyse, le Serveur d'Application prépare les données clients. Nous intégrons un fragment \textit{opt} pour vérifier la santé du Moteur d'Inférence ML via son endpoint dédié.

Si le service est opérationnel, nous initions une communication asynchrone pour transmettre les données par lots. Le Moteur d'Inférence ML charge le modèle Random Forest, exécute les prédictions et calcule les valeurs SHAP. Un fragment \textit{alt} nous permet de gérer les éventuels échecs de connexion ou de calcul : en cas d'erreur du micro-service, nous basculons sur un message d'alerte sans interrompre le serveur principal. Les résultats JSON sont ensuite persistés dans PostgreSQL.

\begin{figure}[htbp]
    \centering
    \includegraphics[width=0.9\linewidth, height=13cm]{images/fastapi.png}
    \caption{Diagramme de séquence de l'analyse sur données réelles }
    \label{fig:ds_analyse_donnees_reelles}
\end{figure}

\subsection{Flux 3~: Workflow de recommandation et validation}

Nous avons modélisé le cycle de vie des actions de rétention pour refléter l'organisation hiérarchique de Tunisie Telecom. Le flux débute par la création d'une recommandation par un agent. Nous enregistrons cette action avec un statut initial de validation requise.

Le système génère alors une notification automatique pour le Chef d'Agence concerné. Ce dernier peut consulter le détail de la proposition et décider de l'activer ou de la rejeter. En cas d'activation, l'action est transmise à l'agent pour exécution. Nous assurons la traçabilité de chaque changement d'état, permettant ainsi un suivi précis de l'effort commercial de l'agence.

\begin{figure}[htbp]
    \centering
    \includegraphics[width=0.9\linewidth, height=9cm]{images/recom.png}
    \caption{Diagramme de séquence du workflow de recommandation}
    \label{fig:workflow_recommandation}
\end{figure}

\subsection{Flux 4~: Exportation de rapports PDF}

Pour faciliter la portabilité des diagnostics, nous avons mis en place un processus de génération de documents à la volée. Lorsque l'utilisateur demande l'export d'une fiche client, notre contrôleur Django récupère l'intégralité du profil client et ses explications SHAP associées depuis PostgreSQL.

Nous injectons ces données dans un template HTML spécialisé que nous transmettons ensuite à la bibliothèque pdfkit. Cette dernière invoque le moteur wkhtmltopdf pour transformer le rendu HTML en un document PDF professionnel. Nous retournons enfin le fichier généré directement au navigateur de l'utilisateur sous forme de flux binaire, assurant un téléchargement instantané sans stockage temporaire sur le serveur.

\begin{figure}[htbp]
    \centering
    \includegraphics[width=0.8\linewidth]{images/export.png}
    \caption{Diagramme de séquence de l'export PDF de la fiche
    client}
    \label{fig:ds_export_pdf}
\end{figure}


\section{Conclusion}

Dans ce chapitre, nous avons présenté la conception complète de notre système. L'architecture hybride Django + FastAPI que nous avons élaborée constitue le choix structurant de notre projet. Nous avons veillé à ce que la modélisation des données reflète fidèlement les contraintes métier. Les diagrammes
que nous avons produits constituent désormais notre référence formelle pour la phase de développement que nous présentons dans le chapitre suivant.
